\section{命题的证明}
\label{app:proofs}

本节给出正文第~\ref{sec:results}~节各命题、推论与注的证明和推导，以及第~\ref{sec:meth-renewal}~节的方差推导。正文的命题、推论、注与式号沿用正文编号，本节新增的式号带前缀 S。

负二项分布 $\NB(r,p)$ 的均值与方差为 $r(1-p)/p$ 与 $r(1-p)/p^2$。取 $r=k$、$p=k/(k+R)$ 即得式~\eqref{eq:offspring}；求和与二项稀疏的结论见第~\ref{sec:meth-model}~节式~\eqref{eq:thinning}，代入 $p=k/(k+R_t)$ 时 $q=(k+\rho R_t)/(k+R_t)$。式~\eqref{eq:relvar} 由式~\eqref{eq:report-dist} 的均值与方差直接相除得到。

\subsection{\texorpdfstring{命题~\ref{prop:current} 的证明}{命题 1 的证明}}

记 $c$ 为 $C_t$ 的观测值，$\theta=\rho R_t$。由式~\eqref{eq:report-dist}，对数似然为 $\ell(R_t)=\text{const}+nk\log[k/(k+\theta)]+c\log[\theta/(k+\theta)]$。对 $R_t$ 求导并整理，得分函数为：
\begin{equation}
  \dot\ell(R_t)=\frac{c}{R_t}-\frac{(c+nk)\rho}{k+\rho R_t}
  =\mathcal I(R_t)\,(\widehat R_t-R_t),\qquad
  \mathcal I(R_t)=\frac{k\rho n}{R_t(k+\rho R_t)}=\frac{n}{R_t/\rho+{R_t}^2/k}.
  \label{eq:score}
\end{equation}

(i) 当 $c>0$ 时，由式~\eqref{eq:score}，得分函数在 $R_t<\widehat R_t$ 时为正，在 $R_t>\widehat R_t$ 时为负，似然在 $\widehat R_t$ 处取唯一最大值。当 $c=0$ 时，$\dot\ell(R_t)=-nk\rho/(k+\rho R_t)<0$，似然在 $(0,\infty)$ 上严格递减，最大值在边界 $R_t=0$，与 $\widehat R_t=0$ 一致。

(ii) 由式~\eqref{eq:report-dist}，$\E\widehat R_t=\rho nR_t/(\rho n)=R_t$，且
\begin{equation}
  \Var(\widehat R_t)=\frac{\rho nR_t+\rho^2nR_t^2/k}{(\rho n)^2}=\frac{R_t}{\rho n}+\frac{R_t^2}{kn}=\frac{R_t(k+\rho R_t)}{k\rho n}.
\end{equation}

(iii) Fisher 信息为得分函数的方差。由式~\eqref{eq:score} 与 (ii) 得
\begin{equation}
  \mathcal I(R_t)=\frac{(k\rho n)^2}{[R_t(k+\rho R_t)]^2}\Var(\widehat R_t)=\frac{k\rho n}{R_t(k+\rho R_t)}=\frac{n}{R_t/\rho+{R_t}^2/k}.
\end{equation}
因此 $\Var(\widehat R_t)=1/\mathcal I(R_t)$。得分函数是 $\widehat R_t-R_t$ 的倍数，这正是 Cramér--Rao 不等式取等号的充要条件\citep{casella2002statistical}，对每个有限 $n$ 都成立。

(iv) 式~\eqref{eq:report-dist} 的概率函数可写为只依赖 $c$ 的因子乘以 $\exp[\upsilon c-\mathcal A(\upsilon)]$，其中自然参数 $\upsilon=\log[\rho R_t/(k+\rho R_t)]$，$\mathcal A$ 为对数配分函数。$k$ 已知时它属于单参数指数族；$R_t$ 取遍 $(0,\infty)$ 时，$\upsilon$ 取遍 $(-\infty,0)$，含开区间，所以 $C_t$ 是完备充分统计量。$\widehat R_t$ 是 $C_t$ 的函数且无偏，由 Lehmann--Scheffé 定理，它是一致最小方差无偏估计\citep{casella2002statistical}。\qed

\subsection{\texorpdfstring{命题~\ref{prop:average} 的证明}{命题 2 的证明}}

由 H4，给定 $n$ 时 $R_t$ 的均值与方差为 $\Rbar$ 与 $v_R$。由命题~\ref{prop:current}(ii)，$\E(\widehat R_t\mid n,R_t)=R_t$，所以 $\E(\widehat R_t\mid n)=\Rbar$。由全方差公式
\begin{equation}
  \Var(\widehat R_t\mid n)=\E\{R_t/(\rho n)+R_t^2/(kn)\mid n\}+\Var(R_t\mid n)
  =\frac{\Rbar}{\rho n}+\frac{\Rbar^2+v_R}{kn}+v_R.
\end{equation}
以 $m=\rho n$ 代入，前两项之和为 $A/m$，得式~\eqref{eq:var-average}。

(i) $A/m\to0$。(ii) 规模项份额为 $(A/m)/(v_R+A/m)=1/(1+mv_R/A)=1/(1+m/\mx)$。(iii) 翻倍降幅为 $[V(m)-V(2m)]/V(m)=[A/(2m)]/(v_R+A/m)=1/[2(1+m/\mx)]$。(iv) $V'(m)=-A/m^2<0$，$V''(m)=2A/m^3>0$；弹性为 $mV'(m)/V(m)=-(A/m)/(v_R+A/m)=-\mx/(m+\mx)$，在 $m=\mx$ 处为 $-1/2$。

(v) 设 $\tilde\theta$ 只依赖本期报告数且 $\E(\tilde\theta\mid R_t)=R_t$。给定 $R_t$，它是 $R_t$ 的无偏估计，由 Cramér--Rao 不等式，$\Var(\tilde\theta\mid R_t)\geq1/\mathcal I(R_t)$。代入式~\eqref{eq:unbiased-bound} 的分解，$\E(\tilde\theta-\Rbar)^2\geq v_R+\E[1/\mathcal I(R_t)]=V(m)$，$\widehat R_t$ 取等号。一步预测时，$y=C_{t+1}/m$，$\hat y$ 由 $\mathcal F_t$ 决定。由 H4，$R_{t+1}$ 与 $\mathcal F_t$ 独立，$C_{t+1}$ 在 $\mathcal F_t$ 下的条件分布与 $\mathcal F_t$ 无关，于是 $\E[(\hat y-y)^2\mid\mathcal F_t]=(\hat y-\E y)^2+\Var(y)\geq\Var(y)=V(m)$，$\hat y=\Rbar$ 取等号。\qed

式~\eqref{eq:pooling} 的推导如下。$J$ 个环境项方差同为 $v_R$、两两相关系数平均为 $\bar c$ 时，其等权平均的方差为 $J^{-2}[Jv_R+J(J-1)\bar cv_R]=v_R[\bar c+(1-\bar c)/J]$。

\subsection{\texorpdfstring{注~\ref{rem:reporting} 的推导}{注 1 的推导}}

把式~\eqref{eq:var-average} 写成 $n$ 与 $\rho$ 的函数
\begin{equation}
  V(n,\rho)=v_R+\frac{\Rbar}{\rho n}+\frac{\Rbar^2+v_R}{kn}.
\end{equation}
令 $\rho\to1$，第二项趋于 $\Rbar/n>0$，第三项与 $\rho$ 无关且为正，故极限严格大于 $v_R$；令 $n\to\infty$，后两项都趋于零。两个方向的边际效应为
\begin{equation}
  -\frac{\partial V}{\partial\log n}=\frac{\Rbar}{\rho n}+\frac{\Rbar^2+v_R}{kn}=\frac{A}{m},\qquad
  -\frac{\partial V}{\partial\log\rho}=\frac{\Rbar}{\rho n}=\frac{\Rbar}{m}.
\end{equation}
二者之比为 $A/\Rbar=1+\rho(\Rbar^2+v_R)/(k\Rbar)$。

\subsection{\texorpdfstring{命题~\ref{prop:allocation} 的证明}{命题 3 的证明}}

(i) 设各期上一代人数均为 $n$，给定 $(R_1,\dots,R_T)$ 时各期报告数相互独立。记 $\bar R_T=T^{-1}\sum_tR_t$。由全方差公式
\begin{equation}
  \frac{1}{T^2}\Var({\textstyle\sum_t}\widehat R_t)=\Var(\bar R_T)+\frac1{T^2}\sum_t\E\{R_t/(\rho n)+R_t^2/(kn)\}
  =\Var(\bar R_T)+\frac{A}{Tm}.
\end{equation}
右端第二项等于 $A/B$，只依赖 $R_t$ 的边际分布，与相关结构无关。$R_t$ 独立时 $\Var(\bar R_T)=v_R/T$，得式~\eqref{eq:alloc-indep}。

一阶自回归情形下，$\Cov(R_s,R_t)=v_R\phi^{|s-t|}$，于是
\begin{equation}
  T^2\Var(\bar R_T)=v_R\,T+2v_R\sum_{h=1}^{T-1}(T-h)\phi^h,\qquad
  \sum_{h=1}^{T-1}(T-h)\phi^h=\frac{\phi[T(1-\phi)-(1-\phi^T)]}{(1-\phi)^2}.
\end{equation}
后一个等式可对 $T$ 归纳验证。代入整理即得式~\eqref{eq:ar-sum}。

单调性。记 $M_T=T^2\Var(\bar R_T)$，$\Delta_j=v_R(1+2\sum_{h=1}^{j}\phi^h)$，则 $M_{T+1}=M_T+\Delta_T$，$M_T=\sum_{j=0}^{T-1}\Delta_j$。$\Var(\bar R_{T+1})<\Var(\bar R_T)$ 等价于 $T^2\Delta_T<(2T+1)M_T$。数列 $\Delta_j$ 的增量 $2v_R\phi^{j+1}$ 递减，所以它是凹数列，满足 $\Delta_j\geq[(T-j)\Delta_0+j\Delta_T]/T$。对 $j=0,\dots,T-1$ 求和，得 $M_T\geq[(T+1)\Delta_0+(T-1)\Delta_T]/2$。记 $d=\Delta_T-\Delta_0$，代入 $\Delta_0=v_R$，有
\begin{equation}
  (2T+1)M_T-T^2\Delta_T\geq(T+1)(Tv_R-d/2).
\end{equation}
由于 $d=2v_R\sum_{h=1}^{T}\phi^h<2Tv_R$，右端为正。$\phi=0$ 时即独立情形，单调性显然成立。

(ii) 由式~\eqref{eq:alloc-indep}，以 $T=\mathcal B/(c_0+c_1m)$、$B=Tm$ 代入，得
\begin{equation}
  f(m)=\frac{(c_0+c_1m)(v_R+A/m)}{\mathcal B}=\frac{c_0v_R+c_1A+c_1v_Rm+c_0A/m}{\mathcal B}.
\end{equation}
$f'(m)=(c_1v_R-c_0A/m^2)/\mathcal B$，$f''(m)=2c_0A/(m^3\mathcal B)>0$。$c_0>0$ 时 $f$ 在 $(0,\infty)$ 上严格凸，唯一驻点为 $m^{*2}=c_0A/(c_1v_R)=(c_0/c_1)\mx$。代回得 $f(m^*)=[c_0v_R+c_1A+2\sqrt{c_0c_1v_RA}]/\mathcal B$，即式~\eqref{eq:optimal-m}。

约束 $T\geq1$ 等价于 $m\leq(\mathcal B-c_0)/c_1$。若 $m^*$ 超出可行范围，$f$ 在可行区间上递减，最优解取右端点，即 $T=1$。若 $c_0=0$，$f(m)=c_1(v_Rm+A)/\mathcal B$ 严格递增，下确界在 $m\to0$ 处取得。

整数解。$m=(\mathcal B/T-c_0)/c_1$ 关于 $T$ 单调，而 $f$ 关于 $m$ 单峰，所以 $f$ 作为 $T$ 的函数也单峰，整数最优解是连续最优期数两侧的两个整数之一，正文的取整例子即按此比较。环境服从一阶自回归时的精确结果见命题~\ref{prop:ar-design}。\qed

\subsection{增长检验的证明}
\label{sec:s-test}

正文第~\ref{sec:res-test}~节的增长检验以 $C_t\geq c_\alpha$ 为拒绝域，下面给出其三条性质的证明。

取 $R>R'$。由式~\eqref{eq:report-dist}，两者的似然比为
\begin{equation}
  \frac{f(c;R)}{f(c;R')}=\frac{(k+\rho R')^{nk}\,[R(k+\rho R')]^{c}}{(k+\rho R)^{nk}\,[R'(k+\rho R)]^{c}}.
\end{equation}
由于 $R/(k+\rho R)$ 关于 $R$ 严格递增，$R(k+\rho R')/[R'(k+\rho R)]>1$，似然比关于 $c$ 严格递增，分布族具有单调似然比。

(a) 水平。对单调似然比族，形如 $\{C_t\geq c\}$ 的检验的功效函数关于 $R_t$ 非降\citep{lehmann2005testing}，所以在 $R_t\leq1$ 上的最大拒绝概率在 $R_t=1$ 处取得，等于 $\alpha^*=P_{R_t=1}(C_t\geq c_\alpha)\leq\alpha$。

(b) 一致最优。由 Karlin--Rubin 定理\citep{karlin1956theory}，对任意水平，形如“$C_t>c$ 时拒绝，$C_t=c$ 时以概率 $\gamma$ 拒绝”的检验在该水平下一致最优。非随机检验 $\{C_t\geq c_\alpha\}$ 是 $c=c_\alpha-1$、$\gamma=1$ 的特例，所以在水平 $\alpha^*$ 下一致最优；适当选取 $\gamma$ 可以使水平恰为 $\alpha$。

(c) 功效。$R_t=1+\delta$ 时，$C_t\sim\NB(nk,k/[k+\rho(1+\delta)])$，拒绝概率为 $1-F_{\NB}(c_\alpha-1)$，即式~\eqref{eq:power}。\qed

\subsection{\texorpdfstring{命题~\ref{prop:report-noise} 与推论~\ref{cor:switch} 的证明}{命题 4 与推论 1 的证明}}

(i) 给定 $(n,R_t,p_t)$，由式~\eqref{eq:report-dist}，$\E(C_t\mid\cdot)=p_tnR_t$，$\Var(C_t\mid\cdot)=p_tnR_t+p_t^2nR_t^2/k$。对 $p_t$ 用全方差公式，并代入 $\E p_t=\rho$、$\E p_t^2=\rho^2+\sigma^2$，得：
\begin{equation}
  \Var(C_t\mid n,R_t)=\rho nR_t+\frac{(\rho^2+\sigma^2)nR_t^2}{k}+\sigma^2n^2R_t^2 .
\end{equation}

(ii) $\E(C_t\mid n,R_t)=\rho nR_t$，所以 $\E\widehat R_t=\Rbar$。再对 $R_t$ 用全方差公式，$\Var(\E(\widehat R_t\mid R_t))=v_R$，条件方差的期望为 $(\rho n)^{-2}\E[\Var(C_t\mid n,R_t)]$。代入 $\E R_t^2=\Rbar^2+v_R$ 与 $m=\rho n$，得：
\begin{equation}
  V_H(m)=v_R+\frac{\sigma^2}{\rho^2}(\Rbar^2+v_R)+\frac{\Rbar}{m}+\frac{(\rho^2+\sigma^2)(\Rbar^2+v_R)}{\rho km},
\end{equation}
整理即为式~\eqref{eq:r3-vH}。第二项不含 $n$。

(iii) 记 $\theta_t=p_tR_t=\rho R_tu_t$。给定 $\theta_t$ 时 $C_t$ 的分布 $\NB(nk,k/(k+\theta_t))$ 与 $R_t$ 无关，$R_t\to\theta_t\to C_t$ 构成马尔可夫链。由 Fisher 信息的数据处理不等式\citep{zamir1998proof}：
\begin{equation}
  \mathcal I_{C_t}(R_t)\ \leq\ \mathcal I_{\theta_t}(R_t)=\frac{\mathcal J_u}{R_t^2},
\end{equation}
其中等号右端是尺度族 $\theta_t/\rho=R_tu_t$ 关于尺度参数的 Fisher 信息，其密度为 $f_u(y/R)/R$，对 $R$ 求导后在 $u=y/R$ 处得分函数为 $-[1+uf_u'(u)/f_u(u)]/R$。对条件无偏估计量用给定 $R_t$ 的 Cramér--Rao 不等式，$\Var(\tilde\theta\mid R_t)\geq R_t^2/\mathcal J_u$；再按式~\eqref{eq:unbiased-bound} 的方式对 $R_t$ 取期望，即得式~\eqref{eq:r3-lb} 的第一个不等式。

记尺度族的得分函数为 $s_u(u)=-[1+uf_u'(u)/f_u(u)]$（取 $R=1$）。$\E u_t=1$ 对尺度参数求导得 $\E[s_u(u_t)\,u_t]=1$，又 $\E s_u(u_t)=0$，故 $\Cov(s_u(u_t),u_t)=1$。由 Cauchy--Schwarz 不等式，$1\leq\mathcal J_u\Var(u_t)=\mathcal J_u\tau^2$，即 $\mathcal J_u\geq1/\tau^2$。等号成立当且仅当 $s_u(u)=b_u(u-1)$，即 $f_u'(u)/f_u(u)=(b_u-1)/u-b_u$，解得 $f_u(u)\propto u^{b_u-1}e^{-b_uu}$，为形状与速率同为 $b_u=1/\tau^2$ 的 Gamma 分布。此时下界为 $v_R+\tau^2(\Rbar^2+v_R)=v_R+v_H$，对每个 $m$ 成立；$\widehat R_t$ 的均方误差为 $v_R+v_H+A_H/m$，$m\to\infty$ 时趋于下界，下确界因此为 $v_R+v_H$。对数正态 $u_t=\exp(sZ-s^2/2)$ 时，$s_u(u)=(\ln u+s^2/2)/s^2$，$\mathcal J_u=1/s^2$，而 $\tau^2=e^{s^2}-1$，故 $(\Rbar^2+v_R)/\mathcal J_u=v_H\log(1+\tau^2)/\tau^2$。


(iv) 给定 $(R_t,p_t)$，$\E(C_t\mid\cdot)=p_tnR_t$，$\tilde R_t=C_t/(p_tn)$ 的条件方差为 $R_t/(p_tn)+R_t^2/(kn)$。条件均值 $R_t$ 的方差为 $v_R$；对 $(R_t,p_t)$ 取期望并利用二者独立，即得式~\eqref{eq:r3-obs}。$\E(1/p_t)$ 有限保证式中第二项有限。二项型波动下 $\E(1/p_t)=B(\rho\nu-1,(1-\rho)\nu)/B(\rho\nu,(1-\rho)\nu)=(\nu-1)/(\rho\nu-1)$，要求 $\rho\nu>1$；$\rho\nu\leq1$ 时 Beta 密度在 0 附近的质量使 $\E(1/p_t)$ 发散。$\rho=0.25$、$\nu=100$ 时 $\E(1/p_t)=4.125$。制度型波动取 $u_t$ 服从形状与速率同为 $1/\tau^2$ 的 Gamma 分布时，$\E(1/u_t)=1/(1-\tau^2)$，要求 $\tau<1$。

(v) 由 (ii)，$v_R+v_H=(1+\tau^2)(\Rbar^2+v_R)-\Rbar^2=U-\Rbar^2$，$A_H=\Rbar+\rho U/k$，$\E\widehat R_t=\Rbar$ 不含 $v_R$ 与 $\tau^2$。$(v_R,\tau^2)$ 沿 $U$ 的等值线变化时，报告数的均值与方差都不变。

(vi) 无计数噪声时 $C_{j,t}=\rho_ju_{j,t}I_{j,t}$，$I_{j,t}$ 为真实感染数。$\rho_j$ 在 $G_t$ 的分子与分母中各出现一次，精确约去。真实感染数在滞后窗口内近似不变时，$G_t/G_t^0=u_{j,t+1}/\sum_sw_su_{j,t+1-s}$。令 $u=1+\epsilon$，由 $\sum_sw_s=1$ 一阶展开得 $G_t/G_t^0-1=\epsilon_{t+1}-\sum_sw_s\epsilon_{t+1-s}+O(\epsilon^2)$。其方差为：
\begin{equation}
  \tau_w^2\Bigl[1-2\sum_sw_s\psi^{s}+\sum_s\sum_{s'}w_sw_{s'}\psi^{|s-s'|}\Bigr]+o(\tau_w^2),
\end{equation}
即式~\eqref{eq:r3-F}。截距比 $q_j=F(w_j,\psi)/F(w_{\text{主}},\psi)$ 只依赖 $\psi$；在 $\psi\in[0,1)$ 的 2001 个格点上求上下确界，得 $q_1\in[0.883,1.338]$，$q_2\in[0.892,1.374]$，$q_2$ 的上确界在 $\psi\to1$ 时取到。

推论~\ref{cor:switch}。由式~\eqref{eq:r3-vH}，把 $\tau^2$ 视为 $\rho$ 的函数：
\begin{equation}
  V_H=v_R+\tau^2(\Rbar^2+v_R)+\frac{\Rbar}{\rho n}+\frac{(1+\tau^2)(\Rbar^2+v_R)}{kn}.
\end{equation}
提高报告率与扩大覆盖的边际效果分别为 $-\rho\,\partial V_H/\partial\rho$ 与 $-n\,\partial V_H/\partial n$，两者之差乘以 $kn/(\Rbar^2+v_R)$ 为：
\begin{equation}
  -\rho\frac{\mathrm d\tau^2}{\mathrm d\rho}(kn+1)-(1+\tau^2),
\end{equation}
$\Rbar/(\rho n)$ 一项在两种边际中相同，已相互抵消。差为正即式~\eqref{eq:r3-switch}。二项型下 $\tau^2=(1-\rho)/[\rho(\nu+1)]$，$-\rho\,\mathrm d\tau^2/\mathrm d\rho=1/[\rho(\nu+1)]$，$1+\tau^2=(\rho\nu+1)/[\rho(\nu+1)]$，条件化为 $kn+1>\rho\nu+1$，即 $kn>\rho\nu$，亦即 $m>\rho^2\nu/k$。制度型下 $\mathrm d\tau^2/\mathrm d\rho=0$，边际之差恒为负；此时 $-n\,\partial V_H/\partial n=\Rbar/m+(1+\tau^2)(\Rbar^2+v_R)/(kn)$，$-\rho\,\partial V_H/\partial\rho=\Rbar/m$，两者之比为 $A_H/\Rbar$。\qed

\subsection{\texorpdfstring{命题~\ref{prop:vR-info} 的有限期数结果与证明}{命题 5 的有限期数结果与证明}}
\label{sec:s-vR}

正文命题~\ref{prop:vR-info}(ii) 只给出 $T\Var(\hat v_R)$ 的极限平台。下面的补充命题给出有限 $T$ 的方差、$\Omega(m)$ 的 $1/m$ 展开与 Fisher 信息上界，记号与正文相同，$\mu_{jY}$ 为 $Y_t=C_t/m$ 的 $j$ 阶中心矩。

\begin{sproposition}[环境方差估计的有限期数精度]
\label{sprop:vR-finite}
在 H1--H4 与 H6 下：
\begin{enumerate}[label=(\alph*)]
  \item 各期 $R_t$ 独立同分布时，$\hat v_R$ 的偏差为 $-\rho\sigma_Y^2/[kmT(1+\rho/(km))]$，方差为
  \begin{equation}
    \Var(\hat v_R)=\frac{1}{(1+\rho/(km))^2}\Bigl[\frac{\mu_{4Y}}{T}-\frac{(T-3)\,\sigma_Y^4}{T(T-1)}-\frac{2\varpi'\mu_{3Y}}{T}+\frac{\varpi'^2\sigma_Y^2}{T}\Bigr]\bigl(1+O(\tfrac{1}{mT})\bigr),
    \label{eq:r3-omega}
  \end{equation}
  其中 $\varpi'=(1+2\rho\Rbar/k)/m$。$R_t$ 服从 Gamma 分布时
  \begin{equation}
    \Omega(m)=\Omega_\infty+\frac{\Omega_1}{m}+O(m^{-2}),\qquad
    \Omega_\infty=v_R^2\Bigl(2+\frac{6v_R}{\Rbar^2}\Bigr),\qquad
    \Omega_1=4v_RA+\frac{8v_R^2\,[k\Rbar+3\rho(\Rbar^2+v_R)]}{k\Rbar^2}.
    \label{eq:r3-c1}
  \end{equation}
  \item 设 $(R_1,\dots,R_T)$ 的联合分布属于以 $v_R$ 为参数、其余参数已知的模型族。对任意规模 $m$ 与任意相关结构，$C_{1:T}$ 关于 $v_R$ 的 Fisher 信息不超过 $R_{1:T}$ 关于 $v_R$ 的 Fisher 信息，$v_R$ 的任一无偏估计满足 $\Var\geq1/\mathcal I_{R_{1:T}}(v_R)$。各期独立且服从均值固定的 Gamma 分布时
  \begin{equation}
    \mathcal I_{R_{1:T}}(v_R)=\frac{T\xi\,[\xi(\ln\Gamma)''(\xi)-1]}{v_R^2},\qquad \xi=\frac{\Rbar^2}{v_R},
    \label{eq:r3-IR}
  \end{equation}
  $(\ln\Gamma)''$ 为三伽马函数；$R_{1:T}$ 服从系数 $\phi$ 已知的平稳高斯一阶自回归时，$\mathcal I_{R_{1:T}}(v_R)=T/(2v_R^2)$，与 $\phi$ 无关。
\end{enumerate}
\end{sproposition}

式~\eqref{eq:r3-c1} 是 $1/m\to0$ 处的一阶展开，只适用于大 $m$。在正文图~\ref{fig:variance} 的参数下 $\Omega_1/\Omega_\infty=108.6$，$m=10$ 时一阶式给出 $\Omega(m)/\Omega_\infty=11.9$，精确值为 37.5；$m=100$ 时为 2.09 与 2.36；$m=1000$ 时精确值为 1.11。$m$ 为几十或更小时应使用式~\eqref{eq:r3-omega} 的精确矩计算。同一参数下每期的 Fisher 信息上界为 $\mathcal I_R=315.6$，单期报告数在 $m=1$、10、100、1000、3000 时的信息为 0.17、8.0、134.9、284.8、304.8，随规模单调上升并趋近上界；相对标准误的下界 $(0.505T)^{-1/2}$ 由此得到。

下面依次证明正文命题~\ref{prop:vR-info} 与补充命题~\ref{sprop:vR-finite} 的各条。
正文命题~\ref{prop:vR-info}(i)。由式~\eqref{eq:var-average}，$\Var(C_t)=m^2V(m)=m^2v_R+m\Rbar+(\rho m/k)(\Rbar^2+v_R)$。把 $\Rbar$ 换成 $\bar C/m$、$\Var(C_t)$ 换成 $S_C^2$ 后对 $v_R$ 求解，得式~\eqref{eq:r3-vhat} 的第一种形式；两边除以 $m^2$ 得第二种形式。

补充命题~\ref{sprop:vR-finite}(a) 与正文命题~\ref{prop:vR-info}(ii)。独立同分布时 $\E S_Y^2=\sigma_Y^2$，$\E\bar Y=\Rbar$，$\E\bar Y^2=\Rbar^2+\sigma_Y^2/T$，所以 $\E\varpi(\bar Y)=\varpi(\Rbar)+\rho\sigma_Y^2/(kmT)$，偏差即由此而来。方差方面，有以下精确结果\citep{casella2002statistical}：
\begin{equation}
  \Var(S_Y^2)=\frac{\mu_{4Y}}{T}-\frac{(T-3)\sigma_Y^4}{T(T-1)},\qquad
  \Cov(S_Y^2,\bar Y)=\frac{\mu_{3Y}}{T},\qquad
  \Var(\bar Y)=\frac{\sigma_Y^2}{T}.
\end{equation}
把 $\varpi(\bar Y)$ 在 $\Rbar$ 处线性化，$\varpi(\bar Y)=\varpi(\Rbar)+\varpi'(\bar Y-\Rbar)+\rho(\bar Y-\Rbar)^2/(km)$，最后一项相对于线性项为 $O(1/(mT))$ 量级，即得式~\eqref{eq:r3-omega}。$Y_t$ 的各阶矩可由负二项阶乘矩精确计算。由概率母函数，$\E[C_t(C_t-1)\cdots(C_t-j+1)\mid R_t]=(nk)(nk+1)\cdots(nk+j-1)\,(\rho R_t/k)^j$，用第二类 Stirling 数换成原点矩后，对 $R_t$ 的分布取期望。$m\to\infty$ 时 $Y_t\to R_t$ 且 $\varpi'\to0$，$\Omega(m)\to\mu_{4R}-v_R^2$。形状为 $\xi$、尺度为 $\vartheta$ 的 Gamma 分布的四阶中心矩为 $3\xi(\xi+2)\vartheta^4$，代入 $\xi=\Rbar^2/v_R$、$\vartheta=v_R/\Rbar$ 得 $\Omega_\infty=v_R^2(2+6v_R/\Rbar^2)$。把各阶矩代入 $\Omega(m)=\mu_{4Y}-\sigma_Y^4-2\varpi'\mu_{3Y}+\varpi'^2\sigma_Y^2$ 并在 $1/m=0$ 处展开到一阶，得式~\eqref{eq:r3-c1} 的 $\Omega_1$，该展开另经符号计算核对。$\Omega(m)$ 的单调性只在 $m\in[1,10^5]$ 上做了数值核对，没有解析证明。

补充命题~\ref{sprop:vR-finite}(b)。$v_R\to R_{1:T}\to C_{1:T}$ 构成马尔可夫链。给定 $R_{1:T}$，由 H1 与 H6，$C_{1:T}$ 的条件分布不含 $v_R$。由数据处理不等式，$\mathcal I_{C_{1:T}}(v_R)\leq\mathcal I_{R_{1:T}}(v_R)$；该不等式不依赖 $m$，也不依赖 $R_{1:T}$ 的相关结构。Gamma 情形下 $(\xi,\vartheta)=(\Rbar^2/v_R,v_R/\Rbar)$，Gamma 族关于 $(\xi,\vartheta)$ 的信息矩阵为 $\begin{pmatrix}(\ln\Gamma)''(\xi)&1/\vartheta\\1/\vartheta&\xi/\vartheta^2\end{pmatrix}$，按链式法则代入 $\mathrm d\xi/\mathrm dv_R=-\Rbar^2/v_R^2$、$\mathrm d\vartheta/\mathrm dv_R=1/\Rbar$，整理得式~\eqref{eq:r3-IR}。高斯情形下 $R_{1:T}-\Rbar\sim\mathcal N(0,v_R\Sigma_\phi)$，$\Sigma_\phi$ 已知，关于 $v_R$ 的信息为 $\tfrac12\mathrm{tr}\{[(v_R\Sigma_\phi)^{-1}\Sigma_\phi]^2\}=T/(2v_R^2)$。

正文命题~\ref{prop:vR-info}(iii)。高斯情形下 $\Cov[(R_s-\Rbar)^2,(R_t-\Rbar)^2]=2v_R^2\phi^{2|s-t|}$。因此：
\begin{equation}
  \Var\Bigl[\frac1T\sum_{t=1}^T(R_t-\Rbar)^2\Bigr]=\frac{2v_R^2}{T^2}\sum_{s=1}^{T}\sum_{t=1}^{T}\phi^{2|s-t|}=\frac{2v_R^2}{T}\,\kappa_T(\phi^2),
\end{equation}
与独立情形的 $2v_R^2/T$ 相比，有效期数为 $T/\kappa_T(\phi^2)$。由 $\kappa_T(x)\to(1+x)/(1-x)$ 得式~\eqref{eq:r3-teff} 的近似。\qed

推论~\ref{cor:season} 的证明如下。$v_S\to(\eta_1,\dots,\eta_S)\to(R_t)\to(C_t)$ 构成马尔可夫链，由数据处理不等式，全部数据关于 $v_S$ 的信息不超过 $\eta_{1:S}$ 的信息 $S/(2v_S^2)$。无偏估计的方差不低于 $2v_S^2/S$，相对标准误不低于 $\sqrt{2/S}$。季均值未知时，样本方差 $(S-1)^{-1}\sum_s(\eta_s-\bar\eta)^2$ 的方差为 $2v_S^2/(S-1)$。\qed

\subsection{\texorpdfstring{命题~\ref{prop:levers} 的证明、一阶自回归形式与数值例}{命题 6 的证明、一阶自回归形式与数值例}}
\label{app:levers-num}\label{sec:s-levers}

$R_t$ 与 $u_t$ 都是一阶自回归、系数分别为 $\phi$ 与 $\psi$ 时，$\gamma_R(h)\gamma_u(h)=(\phi\psi)^h$，利用 $\sum_s\sum_tx^{|s-t|}=T\kappa_T(x)$，式~\eqref{eq:r3-levers} 化为
\begin{equation}
  \E(\bar R_T-\Rbar)^2=\frac1T\bigl[v_R\kappa_T(\phi)+\tau^2\Rbar^2\kappa_T(\psi)+\tau^2v_R\kappa_T(\phi\psi)\bigr]+\frac{A_H}{Tm}.
\end{equation}
取图~\ref{fig:variance} 的参数与 $T=10$、$m=50$、$\tau=0.15$、$\phi=\psi=0.5$ 为基线，均方误差为 $22.5\times10^{-3}$，其中环境项 10.40、报告项 7.85、抽样项 4.29（单位 $10^{-3}$）。规模扩大 10 倍只把总量降到 18.7；期数增加到 4 倍降到 6.2；统一口径降到 14.6；报告偏差持续时期数增加到 4 倍后仍为 33.8，其中报告项 29.8。六种设置的蒙特卡洛（每种 60,000 次）与上式的标准化偏差均小于 1.4（图~\ref{fig:r3-levers}C）。

命题~\ref{prop:levers} 的证明如下。$\bar R_T-\Rbar=T^{-1}\sum_t(C_t/m-u_tR_t)+T^{-1}\sum_t(u_tR_t-\Rbar)$。第一部分在给定 $(R_{1:T},p_{1:T})$ 时均值为零，各期独立（H6），与第二部分不相关，方差为 $T^{-2}\sum_t\E\Var(C_t/m\mid R_t,p_t)$。由命题~\ref{prop:report-noise}(i)(ii) 的计算，$\E\Var(C_t/m\mid R_t,p_t)=A_H/m$，第一部分方差为 $A_H/(Tm)$。第二部分中，$u$ 与 $R$ 独立且 $\E u_t=1$，故：
\begin{equation}
  \Cov(u_sR_s,u_tR_t)=\E(u_su_t)\E(R_sR_t)-\Rbar^2=v_R\gamma_R(h)+\tau^2\gamma_u(h)\bigl[\Rbar^2+v_R\gamma_R(h)\bigr].
\end{equation}
对 $s,t$ 求和即得式~\eqref{eq:r3-levers}。一阶自回归时 $\gamma_R(h)\gamma_u(h)=(\phi\psi)^h$，并利用 $\sum_s\sum_tx^{|s-t|}=T\kappa_T(x)$ 得到简化形式。(i)--(iii) 由式~\eqref{eq:r3-levers} 直接读出。$m$ 只出现在最后一项；$\gamma_u\equiv1$ 时报告项为 $\tau^2[\Rbar^2+v_R\kappa_T(\phi)/T]\geq\tau^2\Rbar^2$；$\tau=0$ 时报告项为零，使用 $p_t$ 的情形见命题~\ref{prop:report-noise}(iv)。(iv) 即命题~\ref{prop:vR-info}。\qed

\subsection{\texorpdfstring{命题~\ref{prop:ar-design} 的证明}{命题 7 的证明}}

由 $Tc_1m=\mathcal B-c_0T$，目标函数可写成：
\begin{equation}
  f_\phi(T)=v_Re(T)+\frac{Ac_1}{\mathcal B-c_0T},\qquad 1\leq T<\mathcal B/c_0 .
\end{equation}

(i) 记 $K(T)=2(1-\phi^T)/T+\phi^T\ln\phi$，则 $g_T=(1+\phi)/(1-\phi)-2\phi K(T)/(1-\phi)^2$。令 $\ell_\phi=-\ln\phi>0$，$x=\ell_\phi T$，直接求导得：
\begin{equation}
  T^2K'(T)=e^{-x}(x^2+2x+2)-2 .
\end{equation}
记 $Q(x)=e^{-x}(x^2+2x+2)$，$Q(0)=2$，$Q'(x)=-x^2e^{-x}<0$，所以 $x>0$ 时 $Q(x)<2$，$K$ 严格递减，$g_T$ 严格递增。$T=1$ 时，$g_1-1=2\phi(\phi-1-\phi\ln\phi)/(1-\phi)^2$；函数 $\phi(1-\ln\phi)$ 的导数为 $-\ln\phi>0$，在 $\phi=1$ 处等于 1，所以 $\phi(1-\ln\phi)<1$，$g_1<1$。$T\to\infty$ 时 $K(T)\to0$，极限为 $(1+\phi)/(1-\phi)$。

凸性方面，$e''(T)=[2g_T-Tg_T']/T^3$。代入 $g_T'=2\phi(2-Q(x))/[(1-\phi)^2T^2]$ 与 $K$ 的表达式，$e''(T)>0$ 等价于：
\begin{equation}
  \frac{1-\phi^2}{\phi}\ >\ \frac{6-e^{-x}(6+4x+x^2)}{T}=\ell_\phi\,W(x),\qquad W(x)=\frac{6-e^{-x}(6+4x+x^2)}{x}.
\end{equation}
令 $z(x)=(6+4x+x^2)e^{-x}-6+2x$，$z(0)=0$，$z'(x)=2-(2+2x+x^2)e^{-x}$。$(2+2x+x^2)e^{-x}$ 的导数为 $-x^2e^{-x}\leq0$，在 $0$ 处等于 2，故 $z'\geq0$，$z\geq0$，即 $W(x)\leq2$，右端不超过 $2\ell_\phi=2\ln\chi$，$\chi=1/\phi>1$。左端等于 $\chi-1/\chi$。函数 $\chi-1/\chi-2\ln\chi$ 在 $\chi=1$ 处为 0，导数为 $(1-1/\chi)^2>0$，所以 $\chi-1/\chi>2\ln\chi$，$e$ 严格凸。

(ii) $Ac_1/(\mathcal B-c_0T)$ 在可行区间上严格凸，与 $v_Re(T)$ 相加后 $f_\phi$ 严格凸，最优解唯一。内点一阶条件为 $v_Re'(T)+Ac_0c_1/(\mathcal B-c_0T)^2=0$。代入 $e'(T)=-g_T/T^2$ 与 $\mathcal B-c_0T=c_1mT$，得 $v_Rg_T=c_0A/(c_1m^2)$，即式~\eqref{eq:r3-mstar}。记 $\Phi(T,\mathcal B)=f_\phi'(T)$，$\partial\Phi/\partial T=f_\phi''>0$，$\partial\Phi/\partial\mathcal B=-2Ac_0c_1/(\mathcal B-c_0T)^3<0$，由隐函数定理 $\mathrm dT^*/\mathrm d\mathcal B>0$。

(iii) $T^\dagger$ 的存在唯一性由 (i) 中 $g_1<1<g_\infty$ 与严格单调性得到。$T^*=T^\dagger$ 时 $m^*=m_0^*$，所需预算为 $T^\dagger(c_0+c_1m_0^*)=\mathcal B^\dagger$。由 $T^*$ 关于 $\mathcal B$ 的单调性，$\mathcal B<\mathcal B^\dagger$ 时 $T^*<T^\dagger$，$g_{T^*}<1$，$m^*>m_0^*$；反之 $m^*<m_0^*$。$\mathcal B\to\infty$ 时 $f_\phi'(1)\to-v_Rg_1<0$，最优解为内点。若此时 $T^*$ 有界，则 $m^*=(\mathcal B/T^*-c_0)/c_1\to\infty$，与 $m^*\leq m_0^*/\sqrt{g_1}$ 矛盾，故 $T^*\to\infty$，$g_{T^*}\to(1+\phi)/(1-\phi)$。

(iv) 严格凸函数在整数点上的最小值位于连续最小点两侧的相邻整数之一。$f_\phi'(1)\geq0$ 时，由凸性 $f_\phi$ 在可行区间上递增，最优解取 $T=1$。\qed

\subsection{\texorpdfstring{注~\ref{rem:spatial} 的推导}{注 2 的推导}}

$J=\mathcal B/(c_0+c_1m)$ 代入式~\eqref{eq:pooling} 加抽样项 $A/(Jm)$，得 $v_R\bar c+[v_R(1-\bar c)+A/m](c_0+c_1m)/\mathcal B$。对 $m$ 求导并令其为零，$c_1v_R(1-\bar c)=c_0A/m^2$，即 $m^*=\sqrt{c_0A/[c_1v_R(1-\bar c)]}$。

\subsection{条件无偏估计类之外的收缩估计}
\label{app:shrink}

估计量对 $R_t$ 条件有偏时，式~\eqref{eq:unbiased-bound} 不再适用。以最简单的线性收缩 $\hat\theta_\lambda=(1-\lambda)C_t/m+\lambda\mu_0$ 为例，$\mu_0$ 为先验中心，$\delta=\mu_0-\Rbar$，记 $V=v_R+A/m$，均方误差及其最小值为：
\begin{equation}
  \mathrm{MSE}(\lambda)=(1-\lambda)^2V+\lambda^2\delta^2,\qquad
  \lambda^*=\frac{V}{V+\delta^2},\qquad
  \mathrm{MSE}(\lambda^*)=\frac{V\delta^2}{V+\delta^2}.
  \label{eq:shrink}
\end{equation}
$\mathrm{MSE}(\lambda^*)<v_R$ 当且仅当 $\delta^2<v_R(1+m/\mx)$，其中 $\mx=A/v_R$。收缩估计对 $R_t$ 的条件偏差为 $\lambda(\mu_0-R_t)$，它以估计目标1上的偏差换取估计目标2上的方差\citep{efron1975data}。

\subsection{更新过程中的方差推导}
\label{sec:s-renewal}

本小节给出正文第~\ref{sec:meth-renewal}~节式~\eqref{eq:info-decomp}、\eqref{eq:mse-decomp} 与 \eqref{eq:b-delta} 的推导。

以下所有期望都在 $\mathcal F_t$ 条件下取，$\Rbar$、$v_R$ 理解为 $R_{t+1}$ 的条件均值与方差。记 $\bar\Lambda=\E(\Lambda_{t+1}\mid\mathcal F_t)$，$m=\rho\bar\Lambda$。第一次使用全方差公式，以 $(\Lambda_{t+1},R_{t+1})$ 为条件
\begin{equation}
  \Var(C_{t+1})=\E(\rho\Lambda_{t+1}R_{t+1}+\rho^2\Lambda_{t+1}R_{t+1}^2/k)+\Var(\rho\Lambda_{t+1}R_{t+1}).
\end{equation}
由条件独立，右端第一项为 $\rho\bar\Lambda\Rbar+\rho^2\bar\Lambda(\Rbar^2+v_R)/k=m\,A$。乘积的方差为
\begin{equation}
  \Var(\Lambda_{t+1}R_{t+1})=\E(\Lambda_{t+1}^2)\E(R_{t+1}^2)-\bar\Lambda^2\Rbar^2
  =\bar\Lambda^2[v_R+\CV^2(\Lambda_{t+1})(\Rbar^2+v_R)].
\end{equation}
两项相加后除以 $m^2$，即得式~\eqref{eq:info-decomp}。

式~\eqref{eq:mse-decomp} 的推导如下。$D_t$ 由 $\mathcal F_t$ 决定，所以 $\E(G_t\mid\mathcal F_t)=\mu_t/D_t$，$\Var(G_t\mid\mathcal F_t)=(\varsigma_t/D_t)^2$。对 $\E[(G_t-\Rbar)^2\mid\mathcal F_t]=\Var(G_t\mid\mathcal F_t)+[\E(G_t\mid\mathcal F_t)-\Rbar]^2$ 再取期望即可。式~\eqref{eq:identity} 是两边同除以 $\Rbar D_t$ 的代数恒等式。

式~\eqref{eq:b-delta} 的推导换用另一个条件。记 $\mathcal H_t$ 为截至第 $t$ 周的真实感染史，并把式~\eqref{eq:dt} 的滞后权重当作感染压力的权重，即 $\Lambda_{t+1}=\sum_s w_sI_{t+1-s}$，$\sum_s w_s=1$。给定 $\mathcal H_t$，$\Lambda_{t+1}$ 已经确定，取 $m=\rho\Lambda_{t+1}$。它与第~\ref{sec:meth-renewal}~节的 $m=\rho\E(\Lambda_{t+1}\mid\mathcal F_t)$ 相差一层报告噪声，式~\eqref{eq:affine} 的拟合又以 $D_t$ 代替 $m$，式~\eqref{eq:b-delta} 描述的正是由此引入的噪声。各周报告独立地服从 $C_{t+1-s}\mid I_{t+1-s}\sim\Bin(I_{t+1-s},\rho)$，所以分母的条件均值与条件方差为
\begin{equation}
  \E(D_t\mid\mathcal H_t)=\rho\Lambda_{t+1},\qquad
  \Var(D_t\mid\mathcal H_t)=\rho(1-\rho)\sum_s {w_s}^2I_{t+1-s}.
\end{equation}
在平稳规模附近取 $I_{t+1-s}\approx\Lambda_{t+1}$，并令 $D_t=m(1+\delta)$，则 $\E\delta=0$，$\E\delta^2=c_w/m$，其中 $c_w=(1-\rho)\sum_s {w_s}^2$，而 $\E\delta^3$ 与 $\E\delta^4$ 都是 $O(m^{-2})$。记 $X=C_{t+1}/m$，于是 $G_t=X/(1+\delta)$。给定 $\mathcal H_t$，$C_{t+1}$ 的报告抽样与此前各周的报告抽样相互独立，$X$ 与 $\delta$ 独立，展开中没有交叉项。与第~\ref{sec:res-average}~节相同，$\E X=\Rbar$，$\Var X=v_R+A/m$。把 $(1+\delta)^{-1}$ 展开到 $\delta^2$ 再取期望，得到
\begin{equation}
  \E(G_t\mid\mathcal H_t)=\Rbar\Bigl(1+\frac{c_w}{m}\Bigr)+O(m^{-2}),\qquad
  \Var(G_t\mid\mathcal H_t)=v_R+\frac{A+(\Rbar^2+3v_R)\,c_w}{m}+O(m^{-2}).
\end{equation}
只保留 $\delta$ 的一阶项时，方差的附加项是 $\E(R_{t+1}^2)\,c_w/m=(\Rbar^2+v_R)\,c_w/m$。$\delta^2$ 的期望与 $1/m$ 同阶，它另外贡献 $2v_Rc_w/m$，并使条件均值上移 $\Rbar c_w/m$。以条件均值的平方归一化，两处二阶效应相互抵消：
\begin{equation}
  \frac{\Var(G_t\mid\mathcal H_t)}{[\E(G_t\mid\mathcal H_t)]^2}
  =\frac{v_R}{\Rbar^2}+\frac{A+(1-\rho)(\Rbar^2+v_R)\sum_s {w_s}^2}{\Rbar^2\,m}+O(m^{-2}).
\end{equation}
上式中 $1/m$ 的系数就是式~\eqref{eq:b-delta}。若按 $L_t$ 的定义以固定的 $\Rbar$ 归一化，系数变为 $[A+(1-\rho)(\Rbar^2+3v_R)\sum_s {w_s}^2]/\Rbar^2$。在第~\ref{sec:res-inflation}~节的情景参数下，主规格因此多出 0.021，约为式~\eqref{eq:b-delta} 的 1.8\%，正文略去了它。

近似成立需要三个条件。权重须归一且规模大致平稳。规模以周增长比 $\zeta$ 变化时，$\sum_s {w_s}^2$ 应换成
\begin{equation}
  \sum_s {w_s}^2\,\frac{I_{t+1-s}}{\Lambda_{t+1}},
\end{equation}
主规格下它在 $\zeta=1.67$ 时为 0.555，在 $\zeta=0.8$ 时为 0.461。各周报告抽样须条件独立；若报告比例存在跨周的共同冲击，$X$ 与 $\delta$ 相关，展开中会多出交叉项。$m$ 还须足够大。对固定的 $m$ 逐点枚举二项分布并截断 $D_t\geq5$，主规格下 $m[\CV^2(G_t\mid\mathcal H_t)-v_R/\Rbar^2]$ 在 $m=20$ 时为 1.223，在 $m=320$ 时为 1.159，按 $1/m$ 外推为 1.155，与式~\eqref{eq:b-delta} 一致；只用当周权重时，$m=20$ 处为 1.972，高阶余项已不可忽略。实证拟合以 $D_t$ 代替 $m$，属于含噪回归元。第~\ref{sec:res-inflation}~节模拟值与预测值的差距来自含噪回归元、$D_t\geq5$ 的截断和小规模处的高阶余项，三者方向不同，部分相互抵消。
